\documentclass[11pt,a4paper]{article}
\usepackage[margin=2.5cm]{geometry}
\usepackage{amsmath,amssymb,amsthm}
\usepackage{booktabs}
\usepackage{microtype}
\usepackage{xcolor}
\definecolor{cA}{HTML}{1F4E8C}
\definecolor{cB}{HTML}{B3261E}
\definecolor{cC}{HTML}{1B7A3D}
\definecolor{cD}{HTML}{B36B00}
\newcommand{\blk}[2]{\texttt{\textbf{\textcolor{#1}{#2}}}}
\newcommand{\slide}[4]{\texttt{\textcolor{black!25}{#1}\textbf{\textcolor{#4}{#2}}\textcolor{black!25}{#3}}}
\usepackage{parskip}
\usepackage[colorlinks=true,linkcolor=blue!50!black,urlcolor=blue!50!black]{hyperref}

\theoremstyle{definition}
\newtheorem{definition}{Definition}
\newtheorem{example}[definition]{Example}
\theoremstyle{plain}
\newtheorem{proposition}[definition]{Proposition}
\newtheorem{theorem}[definition]{Theorem}

\title{Musings}
\author{Robin Langer}
\date{}

\begin{document}
\maketitle
\vspace{-2em}

\section{The area of a circle}

I know that number theory is strictly forbidden. People have wasted their whole
lives on the subject.

I realised that I never once, in my whole university career, calculated the area
of a circle. So I got Claude to do it.

The method is Archimedes. Inscribe a polygon in the circle. Double the number of
sides. Each doubling halves the angle, and the half angle formula turns a square
root into a nested square root. Starting from the octagon and doubling three
times gives the area of the polygon with sixty four sides.
\[
A_{64} = 64 \times \sqrt{\dfrac{1 + \sqrt{\dfrac{1 + \sqrt{\dfrac{1 + \sqrt{\dfrac{1 + \dfrac{\sqrt{2}}{2}}{2}}}{2}}}{2}}}{2}} \times \sqrt{\dfrac{1 - \sqrt{\dfrac{1 + \sqrt{\dfrac{1 + \sqrt{\dfrac{1 + \dfrac{\sqrt{2}}{2}}{2}}}{2}}}{2}}}{2}}
\]
That is $3.136548$, against $\pi = 3.141593$. Keep doubling and it converges.

And this raised a whole number of questions.

\section{Is there a pattern}

The formula is built out of nothing but square roots of whole numbers. So I
wanted to know what a square root actually looks like. The honest answer is its
continued fraction.

Every square root of a non square whole number repeats forever.

\begin{center}
\begin{tabular}{ll}
\toprule
$\sqrt{2}$  & $[\,1;\ \overline{2}\,]$ \\
$\sqrt{3}$  & $[\,1;\ \overline{1,2}\,]$ \\
$\sqrt{5}$  & $[\,2;\ \overline{4}\,]$ \\
$\sqrt{6}$  & $[\,2;\ \overline{2,4}\,]$ \\
$\sqrt{7}$  & $[\,2;\ \overline{1,1,1,4}\,]$ \\
$\sqrt{13}$ & $[\,3;\ \overline{1,1,1,1,6}\,]$ \\
$\sqrt{19}$ & $[\,4;\ \overline{2,1,3,1,2,8}\,]$ \\
$\sqrt{31}$ & $[\,5;\ \overline{1,1,3,5,3,1,1,10}\,]$ \\
$\sqrt{94}$ & $[\,9;\ \overline{1,2,3,1,1,5,1,8,1,5,1,1,3,2,1,18}\,]$ \\
\bottomrule
\end{tabular}
\end{center}

The notation $[a; b, c, d, \ldots]$ means the tower
\[
  a + \cfrac{1}{b + \cfrac{1}{c + \cfrac{1}{d + \cdots}}}
\]
and the bar marks the part that repeats.

Turns out there is a pattern. There are two of them.

\subsection{Every block is a palindrome followed by twice the first term}

Take the repeating block and drop its last term. What is left reads the same
backwards as forwards. The term you dropped is always twice the leading term.
\[
  \sqrt{19} = [\,4;\ \overline{\underbrace{2,1,3,1,2}_{\text{palindrome}},\ 8}\,]
  \qquad
  \sqrt{31} = [\,5;\ \overline{\underbrace{1,1,3,5,3,1,1}_{\text{palindrome}},\ 10}\,]
\]
I had Claude check this for every non square below two thousand. There are no
exceptions. It is a theorem of Galois and Legendre.

\subsection{Numbers near a square have exact formulas}

Five families. In each one $k$ is the whole part of the root. The last two
need $k \geq 2$.

\begin{center}
\begin{tabular}{lll}
\toprule
$n$ & expansion & examples \\
\midrule
$k^2+1$    & $[\,k;\ \overline{2k}\,]$                & $2, 5, 10, 17, 26, 37, 50, 65$ \\
$k^2+2$    & $[\,k;\ \overline{k, 2k}\,]$             & $3, 6, 11, 18, 27, 38, 51$ \\
$k^2+2k$   & $[\,k;\ \overline{1, 2k}\,]$             & $3, 8, 15, 24, 35, 48, 63$ \\
$k^2+k$    & $[\,k;\ \overline{2, 2k}\,]$             & $6, 12, 20, 30, 42, 56, 72$ \\
$k^2+2k-1$ & $[\,k;\ \overline{1, k-1, 1, 2k}\,]$     & $7, 14, 23, 34, 47, 62, 79$ \\
\bottomrule
\end{tabular}
\end{center}

So $\sqrt{50}$ is $[\,7; \overline{14}\,]$ with no calculation at all, because
fifty is seven squared plus one.

\subsection{Where the pattern stops}

The length of the repeating block is not predictable.

\begin{center}
\begin{tabular}{lrrrrrrrrrrr}
\toprule
$n$      & $43$ & $44$ & $45$ & $46$ & $47$ & $48$ & $50$ & $58$ & $61$ & $65$ & $67$ \\
period   & $10$ & $8$  & $6$  & $12$ & $4$  & $2$  & $1$  & $7$  & $11$ & $1$  & $10$ \\
\bottomrule
\end{tabular}
\end{center}

$\sqrt{61}$ needs eleven terms. $\sqrt{65}$ needs one. No formula gives that row.

This is not noise. It is a harder question wearing a disguise.

Cut the tower off after $i$ steps and you get an ordinary fraction $h/k$. Now
compute $h^2 - n k^2$ and watch it as $i$ grows. Here is $\sqrt{13}$.

\begin{center}
\begin{tabular}{rrrrl}
\toprule
$i$ & $a_i$ & $h/k$ & $h^2 - 13k^2$ & \\
\midrule
$0$ & $3$ & $3/1$     & $-4$ & \\
$1$ & $1$ & $4/1$     & $3$  & \\
$2$ & $1$ & $7/2$     & $-3$ & \\
$3$ & $1$ & $11/3$    & $4$  & \\
$4$ & $1$ & $18/5$    & $-1$ & end of the first block \\
$5$ & $6$ & $119/33$  & $4$  & \\
$6$ & $1$ & $137/38$  & $-3$ & \\
$7$ & $1$ & $256/71$  & $3$  & \\
$8$ & $1$ & $393/109$ & $-4$ & \\
$9$ & $1$ & $649/180$ & $1$  & end of the second block \\
\bottomrule
\end{tabular}
\end{center}

The column hits $\pm 1$ at the end of a block and nowhere else. I had Claude
check that for every non square below three thousand and there are no
exceptions. So the repeating block is precisely the stretch you have to walk to
get from one of these to the next.

A block of even length lands on $+1$ at the end of the first block. A block of
odd length lands on $-1$ there, so you must walk a second block to reach $+1$.
\[
  \text{steps to the smallest solution of } x^2 - ny^2 = 1
  \;=\;
  \begin{cases}
    L - 1  & L \text{ even} \\
    2L - 1 & L \text{ odd}
  \end{cases}
\]

So the period tells you how many steps you must take. It does not tell you how
big the answer is. That comes from the terms you walk past, because each step
multiplies the numerator by roughly $a_i$.

\begin{center}
\begin{tabular}{rrl}
\toprule
$n$ & period & smallest $x, y$ \\
\midrule
$2$   & $1$  & $3,\ 2$ \\
$13$  & $5$  & $649,\ 180$ \\
$29$  & $5$  & $9801,\ 1820$ \\
$61$  & $11$ & $1766319049,\ 226153980$ \\
$109$ & $15$ & $158070671986249,\ 15140424455100$ \\
\bottomrule
\end{tabular}
\end{center}

Fermat set $n = 61$ as a challenge to the English in 1657. He picked it because
the block is long and the answer is enormous. Bhaskara the Second had solved it
five hundred years earlier.

Thirteen and twenty nine both have period five, so both need nine steps. The
terms along the way differ, so the answers differ by a factor of fifteen. The
length says how far. The terms say how big.

Greater minds than myself have pondered these questions. But now with Claude a
mediocre mathematician such as myself can consider them.

\section{Sturmian sequences}

Here is the second question. An irrational number is an infinite object. What
does it look like as a sequence of symbols.

Take an infinite sequence of $0$s and $1$s.
\[
  \ell \;=\; 0\,1\,0\,0\,1\,0\,1\,0\,0\,1\,0\,0\,1\,0\,1\,0 \cdots
\]
Slide a window $m$ letters wide along it and write down every distinct block you
see through it. Let $F(m)$ be how many distinct blocks there are.

Do it. Take $m = 3$ and walk the window along the sixteen letters above.

\begin{center}
\begin{tabular}{llll}
\toprule
position & window & block & \\ \midrule
$0$ & \slide{}{010}{0101001001010}{cA} & \blk{cA}{010} & new \\
$1$ & \slide{0}{100}{101001001010}{cB} & \blk{cB}{100} & new \\
$2$ & \slide{01}{001}{01001001010}{cC} & \blk{cC}{001} & new \\
$3$ & \slide{010}{010}{1001001010}{cA} & \blk{cA}{010} & seen \\
$4$ & \slide{0100}{101}{001001010}{cD} & \blk{cD}{101} & new \\
$5$ & \slide{01001}{010}{01001010}{cA} & \blk{cA}{010} & seen \\
$6$ & \slide{010010}{100}{1001010}{cB} & \blk{cB}{100} & seen \\
$7$ & \slide{0100101}{001}{001010}{cC} & \blk{cC}{001} & seen \\
$8$ & \slide{01001010}{010}{01010}{cA} & \blk{cA}{010} & seen \\
\multicolumn{4}{c}{$\cdots$ and nothing new appears after that $\cdots$} \\
\bottomrule
\end{tabular}
\end{center}

Four distinct blocks, so $F(3) = 4$. Four of the eight possible three letter
blocks never appear at all. This sequence is far from arbitrary.

Now widen the window by one letter. For each of the four blocks, collect every
letter seen following it anywhere in the sequence.

\begin{center}
\begin{tabular}{llll}
\toprule
block & occurs at & followed by & becomes \\ \midrule
\blk{cA}{010} & $0,3,5,8,11,13$ & $0$ at $0$, \textbf{1} at $3$ & \blk{cA}{0100} \emph{and} \blk{cA}{0101} \\
\blk{cB}{100} & $1,6,9$         & $1$, every time              & \blk{cB}{1001} \\
\blk{cC}{001} & $2,7,10$        & $0$, every time              & \blk{cC}{0010} \\
\blk{cD}{101} & $4,12$          & $0$, every time              & \blk{cD}{1010} \\
\bottomrule
\end{tabular}
\end{center}

Three of the four pass through unchanged. One does not. \blk{cA}{010} is
followed by a $0$ at position $0$ and by a $1$ at position $3$, so it becomes
two blocks. Say that it \emph{splits}. So $F(4) = 5$.

Here is why the count is worth having. Suppose you are reading $\ell$ one letter
at a time, trying to call the next letter, and you can remember only the last
$m$. Two occurrences of the same block are then indistinguishable to you. If
that block splits, the next letter genuinely differs between the two
occurrences, so you are wrong at one of them. Not through bad judgement. The
letters you can see do not contain the answer.

\begin{definition}
A sequence is \emph{Sturmian} when exactly one block of each length splits.
\[
  F(m) = m + 1 \qquad \text{for every } m \geq 1 .
\]
\end{definition}

So a Sturmian sequence is one where, at every width, there is precisely one
situation your memory cannot resolve. More memory does not remove it. It only
makes it rarer.

\section{One sequence for each irrational}

\begin{theorem}[Morse and Hedlund]
A sequence $\ell$ is Sturmian if and only if there are an irrational
$\alpha \in (0,1)$ and a $\rho \in [0,1)$ with
\[
  \ell_k \;=\; \lfloor (k+1)\alpha + \rho \rfloor - \lfloor k\alpha + \rho \rfloor
  \qquad \text{for all } k \geq 0 .
\]
\end{theorem}

We do not prove this. The slope $\alpha$ is the density of $1$s and is
determined by $\ell$. The intercept $\rho$ only says where you started reading.

So there is exactly one Sturmian sequence for each irrational number, up to a
shift. The sequence above has slope $1/\varphi^2$ and is the Fibonacci sequence,
the fixed point of $0 \mapsto 01$ and $1 \mapsto 0$.

That is a bijection between two things I did not expect to meet. On one side the
irrational numbers. On the other side the symbol sequences that are minimally
unpredictable.

\section{Now take the irrational to be a square root}

This is the question I actually wanted to ask.

The theorem hands me a sequence for every irrational. Section 2 tells me that
square roots are the irrationals whose continued fractions repeat. So what is
special about the sequence when $\alpha = \sqrt{n} - \lfloor \sqrt{n} \rfloor$.

There is a renormalisation here. Take a Sturmian sequence and pass to its
derived sequence, the one recorded only at the sites where the shorter window
could not decide. The derived sequence is Sturmian again, and its slope is
\[
  G(\alpha) \;=\; \left\{ \tfrac{1}{\alpha} \right\} .
\]
This is the Gauss map. On a continued fraction it does exactly one thing. It
deletes the first term.
\[
  [0;\,a_1,a_2,a_3,\ldots] \;\longmapsto\; [0;\,a_2,a_3,a_4,\ldots]
\]
So the tower of derived sequences is the continued fraction, read one term at a
time. For a generic irrational that list never repeats.

For a square root it does.

\begin{center}
\begin{tabular}{lll}
\toprule
slope & orbit under the Gauss map & cycle \\
\midrule
$\sqrt{2}-1$ & $0.414214 \to 0.414214 \to \cdots$ & $1$ \\
$\sqrt{3}-1$ & $0.732051 \to 0.366025 \to 0.732051 \to \cdots$ & $2$ \\
$\sqrt{5}-2$ & $0.236068 \to 0.236068 \to \cdots$ & $1$ \\
$\sqrt{6}-2$ & $0.449490 \to 0.224745 \to 0.449490 \to \cdots$ & $2$ \\
$\sqrt{7}-2$ & $0.645751 \to 0.548584 \to 0.822876 \to 0.215250 \to 0.645751$ & $4$ \\
$\sqrt{8}-2$ & $0.828427 \to 0.207107 \to 0.828427 \to \cdots$ & $2$ \\
\bottomrule
\end{tabular}
\end{center}

The cycle length is the period of the continued fraction. It is the column from
Section 2, the one I said had no formula, now reappearing as the depth at which
a sequence starts repeating itself.

So the square root sequences are the self similar ones. The splitting never
stops, so no finite table of remembered windows predicts them. But the
renormalisation closes after finitely many levels, so a finite recursion does.

That gap is a real place. A periodic sequence is solved by a table. A generic
Sturmian sequence is solved by nothing finite. The square root sequences sit
between the two, and the period of the continued fraction is how far in they
sit.

There is a theorem here which I state and do not prove. It is due to Crisp,
Moran, Pollington and Shallit.

\begin{theorem}
A Sturmian sequence is the fixed point of a finite substitution if and only if
its slope is a quadratic irrational.
\end{theorem}

Which says that the square root sequences are exactly the Sturmian sequences a
finite machine can generate.

\section{What I take from this}

Three things, none of which I would have reached alone.

The first is that the area of a circle is a question about square roots, and a
square root is a repeating pattern with a palindrome in it.

The second is that the length of that pattern is not a curiosity. It is the size
of the smallest solution of a Pell equation, which is why nobody has a formula
for it.

The third is that the same number appears twice, in two settings that have no
business being related. Once as the period of a continued fraction. Once as the
depth at which a symbol sequence becomes self similar.

I do not know whether any of this is new. It is almost certainly not. But I
found it by asking, and the asking took an afternoon rather than a career, which
is the part that seems worth reporting.

\end{document}
